\documentclass[11pt]{article}
\usepackage[margin=1in]{geometry}
\usepackage{amsmath,hyperref}
\title{A constant-source, constant-sink methane trajectory}
\author{Methane surge reconstruction}
\date{}
\begin{document}
\maketitle
The dashed trajectories in He et al. (2026), Fig.~1, motivate a
constant-parameter counterfactual. Holding gross emissions $E$, soil
uptake $L$ (Tg CH$_4$ yr$^{-1}$), and chemical loss frequency $k$
(yr$^{-1}$) fixed gives
\begin{equation}
 \frac{dm_c}{dt}=E-L-km_c,\qquad m_c(0)=m_0.
\end{equation}
For $k>0$, direct integration gives
\begin{equation}
 m_c(t)=m_0+\left(\frac{E-L}{k}-m_0\right)
                  \left[1-\exp(-kt)\right].
\end{equation}
The limit $k=0$ is $m_c(t)=m_0+(E-L)t$. Time is measured in years
from the specified initial burden. No feedback multiplier is introduced:
holding OH fixed specifies the loss frequency explicitly.

The difference $m_c(t)-m_0$ is a concentration-trajectory attribution.
It differs from summing the relaxation term of He et al.'s equation (1),
because that term uses the actual $m_j$, not the counterfactual $m_c$.
Increased emissions also increase subsequent losses, and their effects
on burden decay with chemical lifetime. In continuous time, with constant
$k_0$ and imposed forcing $F(t)$, the perturbation is
\begin{equation}
 m(t)-m_c(t)=\int_0^t\exp[-k_0(t-s)]F(s)\,ds.
\end{equation}
We will report the published annual budget decomposition and the
constant-parameter trajectory separately. Treating them as numerically
identical would introduce an additional approximation.
\end{document}
